\documentclass[12pt,a4paper]{report}
\usepackage[utf8]{inputenc}
\usepackage{amsmath}
\usepackage{amsfonts}
\usepackage{amssymb}
\usepackage{amsthm}
\usepackage{hyperref}

\usepackage{multicol}
\usepackage{fancyhdr}
\usepackage[inline]{enumitem}
\usepackage{tikz}
\usepackage{tikz-cd}
\usetikzlibrary{calc}
\usetikzlibrary{shapes.geometric}
\usepackage[margin=0.5in]{geometry}
\usepackage{xcolor}
\usepackage{listings} 

\hypersetup{
    colorlinks=true,
    linkcolor=blue,
    filecolor=magenta,      
    urlcolor=cyan,
    pdftitle={Tensors},
    pdfpagemode=FullScreen,
    }

%\urlstyle{same}

\newcommand{\CLASSNAME}{Math 5315 -- Numerical Methods for PDE}
\newcommand{\STUDENTNAME}{Paul Carmody}
\newcommand{\ASSIGNMENT}{Midterm Project }
\newcommand{\DUEDATE}{November 12, 2025}
\newcommand{\SEMESTER}{Fall 2025}
\newcommand{\SCHEDULE}{TTh 2:00 to 3:15}
\newcommand{\ROOM}{Crawford 111}

\newcommand{\MMN}{M_{m\times n}}
\newcommand{\FF}{\mathcal{F}}

\pagestyle{fancy}
 	\fancyhf{}
\chead{ \fancyplain{}{\CLASSNAME} }
%\chead{ \fancyplain{}{\STUDENTNAME} }
\rhead{\thepage}
\input{../MyMacros}

\newcommand{\RED}[1]{\textcolor{red}{#1}}
\newcommand{\BLUE}[1]{\textcolor{blue}{#1}}
\definecolor{mygreen}{RGB}{28,172,0} % color values Red, Green, Blue
\definecolor{mylilas}{RGB}{170,55,241}

\begin{document}
\lstset{language=Matlab,%
    %basicstyle=\color{red},
    breaklines=true,%
    morekeywords={matlab2tikz},
    keywordstyle=\color{blue},%
    morekeywords=[2]{1}, keywordstyle=[2]{\color{black}},
    identifierstyle=\color{black},%
    stringstyle=\color{mylilas},
    commentstyle=\color{mygreen},%
    showstringspaces=false,%without this there will be a symbol in the places where there is a space
    numbers=left,%
    numberstyle={\tiny \color{black}},% size of the numbers
    numbersep=9pt, % this defines how far the numbers are from the text
    emph=[1]{for,end,break},emphstyle=[1]\color{red}, %some words to emphasise
    %emph=[2]{word1,word2}, emphstyle=[2]{style},    
}

\begin{center}
	\Large{\CLASSNAME -- \SEMESTER} \\
	\large{ w/Professor Du}
\end{center}
\begin{center}
	\STUDENTNAME \\
	\ASSIGNMENT -- \DUEDATE\\
\end{center} 
\HLINE

\textbf{Question1.} (40 points)

Apply the FTFS scheme to the wave equation $u_t-u_x=0$, with periodic boundary condition over domain $[0,1]$.  Choose $\Delta t = 0.008$ and $\Delta x=0.01$.
\begin{enumerate}[label=(\alph*)]

	\item Clearly plot the dissipation and dispersion curves of the scheme and list the Matlab code.
	
	\item For the initial condition $u(x,0) = \sin^{40}(\pi x)$, plot the analytical and numerical solution at time $t=1$, quantify the amount of dissipation.
	
	\item Change the initial condition to $u(x,0)=1$ for $0.4\le x \le 0.6$ and $u(x,0)=0$ otherwise.  Plot the analytical and numerical solution at time $t=1$, quantify the amount of dissipation.
	
	\item Compare the results in (b) and (c) and explain the difference based on the Power Spectrum Density Analysis of the initial conditions.

\end{enumerate}

\textbf{Question2.} (40 points)

Apply the LeapFrog scheme ot the wave equation $u_t+u_x=0$, with periodic boundary condition over domain $[0,1]$.  The initial condition is given by $u(x,0)=\sin^{80}(\pi x)$.  Choose $\Delta t = 0.008$ and $\Delta x=0.01$.

\begin{enumerate}[label=(\alph*)]

	\item Analyze the dissipation and dispersion properties for the scheme.
	
	\item Plot the analytical and numerical solution at $t=0.8, t=10,$ and $t=20$.
	
	\item Explain the numerical solution base on the analysis in (a).
	
	\item Can you find a way to reduce the error in the numerical solution?  Confirm your idea with numerical experiments.

\end{enumerate}
\end{document}